\documentclass[10pt]{article}
\usepackage[margin=1in]{geometry}
\usepackage{amsmath,amsthm,amssymb}
\usepackage{mathtools,bbm}
\usepackage{graphicx}
\usepackage{subcaption}
\usepackage{algorithmicx}
\usepackage{algorithm}
\usepackage{algpseudocode}
\usepackage[colorlinks,linkcolor=blue]{hyperref}
\usepackage[noabbrev]{cleveref}
\usepackage{courier}
\usepackage{enumitem}
\usepackage[parfill]{parskip}
\usepackage{multicol}
\usepackage{listings}
\usepackage{color}
\usepackage{xcolor}



\setlist[itemize]{itemsep=.01cm, topsep=.01cm}
\setlist[enumerate]{itemsep=.01cm, topsep=.01cm}


\oddsidemargin 0in
\evensidemargin 0in
\textwidth 6.5in
\topmargin -0.5in
\textheight 9.0in
\setlength{\parskip}{8pt}

\newcommand{\head}[4]{
\thispagestyle{plain}
\newpage
\setcounter{page}{1}
\noindent
\begin{center}
\framebox{ \vbox{ \hbox to 6.28in
{CIS 4190/5190: Applied Machine Learning \hfill #1}
\vspace{4mm}
\hbox to 6.28in{\hfill \large\bf #2 \hfill}% % centered title
\vspace{4mm}
\hbox to 6.28in
{{\it #3 \hfill #4}}
}}
\end{center}
}


\begin{document}
\head{Spring 2026}{News Source Classification from Headlines}{Team Members: Isaac Dcruz, Arjun Verma, Maxwell Zhang}{Project: News Source}

\section{Introduction}
In this project, we built an end-to-end binary NLP classification pipeline to predict whether a headline came from Fox News or NBC News. Our core objective was to make the system reproducible and extensible: every major stage (URL collection, scraping, preprocessing, splitting, training, evaluation, and prediction) is implemented as a script with explicit inputs/outputs and logs.

We first reproduced the handout baseline (\texttt{TF-IDF(max\_features=100) + LogisticRegression(max\_iter=100)}), then expanded to a controlled experiment framework covering multiple feature/model families. Our training framework evaluates a large configuration space and selects the best model by validation macro-F1 before retraining on train+validation and evaluating once on the held-out test set.

This report reflects the current repository state and includes finalized implementation details and experimental outcomes produced in our latest run cycle.

\section{Core Components}

\subsection{Data Collection and Dataset}
Our data pipeline follows a two-stage collection process designed to maximize coverage while keeping source labels reliable. First, we discover article URLs using \texttt{src/collect\_urls.py}. This script traverses robots-discovered sitemap indexes, configured sitemap seeds, RSS/Atom feeds, and section pages for FoxNews and NBC domains. We use this broad discovery strategy because a single source (for example, only one sitemap) misses many valid article pages, especially for historical coverage.

Second, we scrape headline text with \texttt{src/scrape.py}. The scraper fetches each URL with retry/backoff logic and then applies a robust extraction cascade so that template variation across pages does not cause widespread failures:
\begin{enumerate}
    \item source-specific \texttt{h1} class heuristics,
    \item \texttt{itemprop=headline},
    \item first generic \texttt{h1},
    \item \texttt{og:title},
    \item fallback \texttt{title}.
\end{enumerate}

At the raw-data stage, our repository contains:
\begin{itemize}
    \item \texttt{data/raw/original\_urls.csv}: 172,488 rows (115,266 FoxNews; 57,222 NBC).
    \item \texttt{data/scraped/raw\_scraped\_headlines.csv}: 84,825 rows (82,126 success; 2,699 fetch failures), currently only FoxNews labels in file.
    \item \texttt{data/scraped/raw\_scraped\_headlines\_nbc.csv}: 47,446 rows (47,205 success; 241 fetch failures), NBC-only.
\end{itemize}

We then consolidate and preprocess using \texttt{src/preprocess.py}. The preprocessing stage enforces quality and leakage controls by applying scrape-status filtering, null/empty removal, HTML unescaping/tag stripping, whitespace normalization, within-source deduplication, cross-source duplicate removal, and label validation. In addition to the canonical cleaned text, we generate multiple preprocessing variants to support fair ablation studies:
\texttt{headline\_minimal}, \texttt{headline\_lowercase}, \texttt{headline\_nopunct}, \texttt{headline\_nostop}, \texttt{headline\_lemma}.

For internal evaluation, we use \texttt{src/split.py} to produce stratified 70/15/15 train/val/test splits with fixed seed 42 and store split metadata for reproducibility.

After merging Fox and NBC scraped files into \texttt{raw\_scraped\_headlines\_merged.csv}, the final processed dataset contains 128,861 headlines with stratified splits (train: 90,202; val: 19,329; test: 19,330). We also identified 7,940 very short headlines (\textless 3 words). Instead of dropping these automatically, we retain them and account for their behavior in error analysis, since short headlines are a meaningful real-world edge case.

\subsection{Model Design}
We designed model development as a staged but fully scripted process from baseline reproduction to representation-rich and ensemble architectures. All model files in \texttt{src/models} share the same training contract through \texttt{src/models/\_base.py}: each \texttt{ModelConfig} specifies \texttt{name}, \texttt{prep\_col}, and an unfitted sklearn estimator; the runner trains on \texttt{train.csv}, ranks by validation macro-F1, retrains the selected configuration on train+val, and evaluates exactly once on held-out test. This standardized protocol ensures direct comparability across families.

\textbf{Global training/evaluation protocol (applies to every classical family):}
\begin{itemize}
    \item Fixed label encoding (\texttt{FoxNews=0}, \texttt{NBC=1}) and fixed split seed (42).
    \item Selection metric: validation macro-F1 (\texttt{val\_f1\_macro}).
    \item Final reporting metrics: accuracy, macro/weighted F1, per-class F1, macro precision/recall, confusion matrix.
    \item Refit policy: winner retrained on train+val before single test evaluation.
    \item Artifact policy: save model (\texttt{models/*\_best.joblib}) and metadata JSON with winning config and test metrics.
\end{itemize}

\textbf{Baseline model (\texttt{src/train\_baseline.py}):}
\begin{itemize}
    \item Handout-consistent baseline: \texttt{TfidfVectorizer(max\_features=100)} + \texttt{LogisticRegression(max\_iter=100)}.
    \item Used as a low-capacity anchor to quantify absolute gains from richer representations and stronger learners.
\end{itemize}

\textbf{Classical model families (all implemented and executed):}
\begin{enumerate}
    \item \textbf{Word TF-IDF + Logistic Regression (\texttt{tfidf\_logreg.py}, 25 configs).}
    We sweep vocabulary size (\texttt{max\_features} from 100 to 20,000), word n-gram ranges \((1,1)\), \((1,2)\), \((1,3)\), \((2,2)\), regularization strength \(C\in\{0.01,\dots,20\}\), and preprocessing variants (\texttt{minimal/lowercase/nopunct/nostop/lemma}). All variants use \texttt{sublinear\_tf=True}; stopword handling is adjusted to avoid confounding in already de-stopped columns.

    \item \textbf{Word TF-IDF + LinearSVC (\texttt{tfidf\_svm.py}, 22 configs).}
    We evaluate large-margin linear SVMs on sparse TF-IDF with sweeps over \(C\), \texttt{max\_features}, and n-gram ranges. Because raw \texttt{LinearSVC} has no probability output, calibrated variants (\texttt{CalibratedClassifierCV}, sigmoid/Platt scaling, CV=3) are included for downstream soft-voting/stacking compatibility.

    \item \textbf{Word TF-IDF + MultinomialNB (\texttt{tfidf\_nb.py}, 22 configs).}
    We sweep smoothing \(\alpha\), feature counts, n-gram ranges, and preprocessing variants. Critically, all NB pipelines set \texttt{norm=None} in TF-IDF to preserve count-like magnitude semantics required by multinomial likelihood modeling.

    \item \textbf{Character n-gram models (\texttt{char\_ngram.py}, 24 configs).}
    We use \texttt{analyzer=char\_wb} (token-boundary-aware character n-grams) with LR/SVM, ranges such as \((2,4)\) to \((5,7)\), and feature-count sweeps up to 20,000. We explicitly test capitalization-preserving input (\texttt{headline\_minimal}) vs lowercased input, and include combined word+char \texttt{FeatureUnion} pipelines to exploit complementary lexical and orthographic cues.

    \item \textbf{Stylometric tabular models (\texttt{stylometric.py}, 12 configs).}
    We extract 15 handcrafted headline-style features (e.g., length statistics, punctuation frequencies, capitalization ratios, digit ratio, title-case ratio) and train LR, LinearSVC, RandomForest, and HistGradientBoosting variants. This family isolates writing-style signal independent of specific vocabulary identity.

    \item \textbf{Hybrid sparse unions (\texttt{hybrid.py}, 10 configs).}
    We combine word TF-IDF, optional char TF-IDF, and stylometric features through \texttt{FeatureUnion}, using \texttt{SparseWrapper} to keep stylometric outputs CSR-compatible. A \texttt{MaxAbsScaler} is applied before the classifier so low-magnitude style features are not drowned out by high-dimensional TF-IDF coordinates.

    \item \textbf{Soft-voting ensembles (\texttt{voting\_ensemble.py}, 7 configs).}
    We average class probabilities from diverse base learners (word LR, char LR, MNB, calibrated SVM, stylometric LR), including weighted voting (\([2,1,1]\)) and feature-diverse branches. Minimal-input variants preserve capitalization for stylometric usefulness while normalizing TF-IDF branches as needed.

    \item \textbf{Stacking ensembles (\texttt{stacking\_ensemble.py}, 7 configs).}
    We train meta-learners on out-of-fold base-model probabilities (\texttt{stack\_method='predict\_proba'}) with CV=5 (or CV=3 for the heaviest 5-base stack). Meta-learners include LR and HistGradientBoosting; base learners include LR, char models, MNB, calibrated SVM, and stylometric branches depending on configuration.
\end{enumerate}

Across these eight classical families we execute \(25+22+22+24+12+10+7+7=129\) configurations, matching the full classical sweep reported in this project cycle.

\textbf{Transformer baselines (\texttt{src/train\_transformer.py}):}
\begin{itemize}
    \item DistilBERT (\texttt{distilbert-base-uncased}) and DeBERTa-v3-small (\texttt{microsoft/deberta-v3-small}) are trained with Hugging Face \texttt{Trainer}.
    \item Training script includes deterministic seeding, tokenizer fallback (\texttt{use\_fast=True} \(\rightarrow\) \texttt{use\_fast=False}), padded batching, epoch-level evaluation, best-checkpoint loading by macro-F1, and saved confusion-matrix artifacts.
    \item These runs provide a compute-aware SOTA reference against the strongest classical pipeline on the same task formulation.
\end{itemize}

\subsection{Evaluation and Model Performance}
Our evaluation stack is implemented at three levels. First, \texttt{train\_baseline.py} provides exact baseline reproduction metrics and feature inspection. Second, \texttt{train\_best\_model.py} and family-specific runners produce full experiment tables, selection metadata, and best-model metrics/plots under a common selection protocol. Third, \texttt{evaluate.py} supports deeper post-hoc diagnostics including misclassified samples, confidence-ranked errors, and headline-length error trends.

Concretely, the tooling exports:
\begin{itemize}
    \item scalar metrics (accuracy, macro-F1, per-class F1, precision/recall),
    \item confusion matrices,
    \item ROC AUC and ROC curves from probability/score outputs,
    \item model comparison plots,
    \item feature analyses where supported by model type,
    \item error-analysis artifacts for qualitative review.
\end{itemize}

\textbf{Final reported metrics:}
\begin{itemize}
    \item Baseline test metrics (completed): accuracy = 0.7051, macro-F1 = 0.6481, FoxNews F1 = 0.7897, NBC F1 = 0.5066.
    \item Best-model metrics (completed): model = \texttt{HYB\_word12\_char36\_mf5000\_style\_LR}, validation macro-F1 = 0.8296, test accuracy = 0.8420, test macro-F1 = 0.8294, FoxNews F1 = 0.8758, NBC F1 = 0.7829.
    \item Model comparison figures (completed): \texttt{baseline\_confusion\_matrix.png}, \texttt{best\_confusion\_matrix.png}, \texttt{experiment\_comparison.png}, and \texttt{baseline\_vs\_best.png}.
\end{itemize}

Relative to baseline, our best classical model improves test accuracy by approximately 13.7 percentage points (0.7051 $\rightarrow$ 0.8420) and macro-F1 by approximately 18.1 points (0.6481 $\rightarrow$ 0.8294). This result is substantial and consistent with our ablation logic: combining word n-grams, character n-grams, and stylometric signals captures complementary information that the minimal baseline cannot represent.

To address full model-family coverage, we additionally executed every standalone training entry point under \texttt{src/models} (\texttt{tfidf\_logreg.py}, \texttt{tfidf\_svm.py}, \texttt{tfidf\_nb.py}, \texttt{char\_ngram.py}, \texttt{stylometric.py}, \texttt{hybrid.py}, \texttt{voting\_ensemble.py}, \texttt{stacking\_ensemble.py}) and confirmed all eight completed successfully in \texttt{reports/model\_family\_run\_summary.json}. For each resulting model artifact, we generated a dedicated folder under \texttt{reports/model\_breakdown/<model\_name>/} containing source metrics, recomputed full-test metrics, per-class report JSON, confusion matrix, ROC curve, and copied related figures.

\begin{table}[h]
\centering
\small
\begin{tabular}{lccc}
\hline
Model & Accuracy & Macro-F1 & ROC AUC \\
\hline
baseline & 0.7051 & 0.6481 & 0.7194 \\
tfidf\_logreg\_best & 0.8014 & 0.7812 & 0.8779 \\
tfidf\_svm\_best & 0.8133 & 0.7973 & 0.8870 \\
tfidf\_nb\_best & 0.7740 & 0.7669 & 0.8530 \\
char\_ngram\_best & 0.8303 & 0.8146 & 0.9079 \\
stylometric\_best & 0.7457 & 0.7295 & 0.8229 \\
hybrid\_best & 0.8420 & 0.8294 & 0.9202 \\
voting\_ensemble\_best & 0.8124 & 0.8014 & 0.8909 \\
stacking\_ensemble\_best & 0.8282 & 0.8162 & 0.9085 \\
distilbert & 0.8058 & 0.7948 & 0.8862 \\
deberta\_v3\_small & 0.8024 & 0.7931 & 0.8869 \\
\hline
\end{tabular}
\caption{Held-out test performance for all model families and transformer baselines with generated per-model breakdown folders.}
\end{table}

\subsection{News Source Track Requirement Compliance}
The project description for the News Source track requires a complete ML pipeline (data collection/cleaning/splitting, baseline reproduction and improvement, rigorous evaluation, at least one exploratory component, and leaderboard-oriented model selection). Our implementation and report satisfy these requirements as follows:
\begin{itemize}
    \item \textbf{Baseline reproduction and improvement requirement:} We reproduced the handout baseline protocol (\texttt{TF-IDF max\_features=100 + LogisticRegression max\_iter=100}), and our baseline run achieves test accuracy 0.7051, which exceeds the handout reference of approximately 0.6649. We then improved substantially beyond baseline with multiple model families and ensembles.
    \item \textbf{Data collection and cleaning requirement:} We scraped headline text from the provided source URLs and expanded coverage with additional collection. We documented filtering, deduplication, normalization, and split creation in the data section, and we provide final dataset statistics and split sizes.
    \item \textbf{Evaluation protocol requirement:} We use a fixed split protocol and validation macro-F1 for model selection, followed by a single held-out test evaluation after retraining on train+val, preventing implicit test leakage.
    \item \textbf{Exploratory component requirement:} We include multiple exploratory directions (continuous refresh pipeline, broad architecture ablations, backfill strategy, error analysis, and transformer/SOTA comparisons).
    \item \textbf{Line-chart requirement for report:} The required model-metric trend visualizations (including baseline reference) are provided in \texttt{reports/figures/experiment\_comparison.png} and \texttt{reports/figures/baseline\_vs\_best.png}. The experiment chart includes a baseline reference level and model-family comparisons.
\end{itemize}

\subsection{Leaderboard Submission Record (Required)}
The project handout requires at least one leaderboard submission for News Source Classification. We select the submission model using validation macro-F1 from internal experiments and then generate predictions with \texttt{src/predict.py}. Fill the finalized submission metadata below before final turn-in:
\begin{itemize}
    \item \textbf{Submission model artifact:} \texttt{models/best\_model.joblib} (\texttt{HYB\_word12\_char36\_mf5000\_style\_LR}).
    \item \textbf{Prediction generation command:} \texttt{python src/predict.py --model models/best\_model.joblib --input <leaderboard\_input.csv> --output <predictions.csv>}.
    \item \textbf{Leaderboard submission ID:} \emph{[fill before submission]}.
    \item \textbf{Submission date/time:} \emph{[fill before submission]}.
    \item \textbf{Leaderboard score/metric:} \emph{[fill before submission]}.
\end{itemize}

%\subsection{Model Design} 

%\subsection{Evaluation and Model Performance} 

\section{Exploratory Components}

Beyond the baseline requirement, we explored several complementary directions to understand the task more deeply and improve robustness:
\begin{itemize}
    \item \textbf{Code/System Exploration}: We built a continuous refresh pipeline (\texttt{continuous\_scrape.py}) that repeatedly collects new URLs, merges unseen records, resumes scraping, and logs cycle-level behavior. This makes the dataset growth process repeatable instead of ad hoc.
    \item \textbf{Technique Exploration}: We implemented broad ablations across word n-grams, char n-grams, stylometric features, hybrid representations, voting, and stacking. This allowed us to identify which representational signals and model combinations actually contribute to macro-F1 gains.
    \item \textbf{Dataset Exploration}: We added backfill collection with per-source cursor state so we can progressively expand historical coverage and avoid repeatedly scraping the same recency window.
    \item \textbf{Analysis Exploration}: We added error-analysis artifacts (misclassified samples, confidence-ranked errors, and headline-length error trends) to diagnose failure modes beyond a single aggregate score.
    \item \textbf{SOTA-Oriented Extension}: We trained two transformer baselines with fixed-seed sampled splits on CPU to keep runtime practical in this run cycle.
    \begin{itemize}
        \item DistilBERT (\texttt{distilbert-base-uncased}, 10k/3k/3k train/val/test samples, 1 epoch): test accuracy = 0.8010, test macro-F1 = 0.7908, training runtime $\approx$ 610s.
        \item DeBERTa-v3-small (\texttt{microsoft/deberta-v3-small}, 5k/2k/2k train/val/test samples, 1 epoch): test accuracy = 0.8040, test macro-F1 = 0.7956, training runtime $\approx$ 619s.
        \item In this constrained CPU setting, both transformer runs underperform the best classical hybrid model (macro-F1 0.8294 on full test split). However, they remain meaningful baselines, and a full-data multi-epoch GPU run is a natural next step before making strong architecture-level conclusions.
    \end{itemize}
\end{itemize}

\section{Team Contributions}
\begin{itemize}
    \item \textbf{Isaac Dcruz}: \emph{[Replace with finalized contribution statement.]}
    \item \textbf{Arjun Verma}: \emph{[Replace with finalized contribution statement.]}
    \item \textbf{Maxwell Zhang}: Integrated the ML pipeline components in this repository and led the experimental/evaluation/reporting workflow integration.
\end{itemize}

\section{Execution Summary and Remaining Items}
We finalized this report using one canonical data pipeline, one classical best-model selection flow, and two transformer baselines for SOTA comparison.

\subsection*{Phase 1 --- Data Finalization}
\begin{itemize}
    \item[$\checkmark$] Merge Fox and NBC scraped files into one canonical raw scraped dataset.
    \item[$\checkmark$] Run preprocessing end-to-end to generate \texttt{clean\_headlines.csv}, all text variants, and \texttt{dataset\_summary.txt}.
    \item[$\checkmark$] Regenerate stratified train/val/test splits and verify class balance in \texttt{split\_metadata.json}.
\end{itemize}

\subsection*{Phase 2 --- Classical Model Track}
\begin{itemize}
    \item[$\checkmark$] Run \texttt{train\_baseline.py} and save baseline artifacts.
    \item[$\checkmark$] Run \texttt{train\_best\_model.py} (full, not \texttt{--fast}) to execute all classical configs (129 after removing two incompatible sparse+HGB hybrid configs).
    \item[$\checkmark$] Run every standalone model-family script in \texttt{src/models} and persist per-family best artifacts/metadata.
    \item[$\checkmark$] Verify standalone run completion status: 8/8 scripts successful (\texttt{reports/model\_family\_run\_summary.json}).
    \item[$\checkmark$] Run \texttt{evaluate.py} for baseline and best model to produce full diagnostics and figures.
    \item[$\checkmark$] Lock final classical winner based on validation macro-F1 and held-out test performance.
\end{itemize}

\subsection*{Phase 3 --- SOTA Comparison Track}
\begin{itemize}
    \item[$\checkmark$] Train \textbf{DistilBERT} baseline on the same split protocol (sampled run).
    \item[$\checkmark$] Train \textbf{DeBERTa-v3-small} baseline on the same split protocol (sampled run).
    \item[$\checkmark$] Compare transformer models vs.\ best classical model using the same metrics (accuracy, macro-F1, per-class F1, confusion matrix).
    \item[$\checkmark$] Report compute/performance tradeoff (training time and practical deployment considerations).
\end{itemize}

\subsection*{Phase 4 --- Report Packaging}
\begin{itemize}
    \item[$\checkmark$] Fill all pending metric placeholders and figure references in this document.
    \item[$\checkmark$] Generate per-model breakdown folders (\texttt{reports/model\_breakdown}) with metrics JSON, confusion matrix, and ROC curve for each model artifact.
    \item[$\square$] Embed at least one line-chart figure of model metrics in the final PDF (including baseline), e.g., \texttt{experiment\_comparison.png} and/or \texttt{baseline\_vs\_best.png}.
    \item[$\square$] Confirm final report length/format compliance (5 pages, required course template/format).
    \item[$\square$] Replace team contribution placeholders with finalized one-sentence contributions.
    \item[$\square$] Add at least one leaderboard submission record and evidence (submission ID/date/score and model config used).
    \item[$\square$] Compile final PDF and run a final reproducibility sanity check on all commands listed in the report.
\end{itemize}

\subsection*{Phase 5 --- Final Deliverables Checklist}
\begin{itemize}
    \item[$\square$] Upload the final dataset used for experiments.
    \item[$\square$] Upload the best-performing model artifact/checkpoint.
    \item[$\square$] Upload the final project report PDF.
\end{itemize}

\section{Extra Credit}
Optional demo video (max 10 minutes, all team members present):\\
\textbf{Link:} \emph{[Add final video URL here if submitted.]}

\end{document}